\documentclass[a4paper]{article}
\usepackage{amsfonts}
\usepackage{amsmath}
\usepackage{amssymb}
\usepackage{graphicx}
\newcommand{\dint}{\displaystyle\int}
\newcommand{\limit}{\lim\limits}

\begin{document}
  
\noindent \large Auteur: Tala Mohanna \\
\noindent \large Studierichting: Informatica\\
\noindent \large Oefeningenreeks 6G nr. 1.1\\

\medskip

\normalsize

$ f_n = \mathbf{1}_{\left[0, \tfrac{1}{n}\right]} $ \\

\begin{enumerate}

\item Kwadratisch gemiddelde\\

$ \dint |f_n(x) - 0|^2 dx = \dint_0^{1/n} 1^2 dx $ \\

$ = \left[ x \right]_0^{1/n} = \dfrac{1}{n} $  \\

$ \lim\limits_{n \rightarrow \infty} \dfrac{1}{n} = 0 $ \\

$ \Rightarrow $ deze rij in kwadratisch gemiddelde convergeert naar de nulfuctie. \\

\item Puntsgewijze convergentie\\

Neem $ x = 0 $: \\

$ f_n(0) = 1 $ voor alle $ n $ \\

$ \lim\limits_{n \rightarrow \infty} f_n(0) = 1 \neq 0 $ \\

$ \Rightarrow (f_n)_n $ convergeert naar de nulfunctie maar niet puntsgewijs. \\

\end{enumerate}

\begin{thebibliography}{999}
%\bibitem{a} Referentie
\end{thebibliography}

\end{document}